% Luc Destombe --- licence CC BY-NC-SA 4.0
% https://creativecommons.org/licenses/by-nc-sa/4.0/deed.fr
% Fichier autonome : compiler avec pdflatex, deux fois (signets, nombre de pages).
\documentclass[11pt,a4paper]{article}
\makeatletter
\newif\ifeleve
\elevefalse
\elevetrue

\newif\iflycee@palatino
\lycee@palatinotrue

\RequirePackage[utf8]{inputenc}
\RequirePackage[T1]{fontenc}
\RequirePackage[french]{babel}
\RequirePackage{etoolbox}

\newcommand{\ShorthandsOff}{\@ifundefined{shorthandoff}{}{\shorthandoff{;:!?}}}
\newcommand{\ShorthandsOn}{\@ifundefined{shorthandon}{}{\shorthandon{;:!?}}}
\ShorthandsOff
\AtBeginDocument{\ShorthandsOn}
\AtBeginEnvironment{tikzpicture}{\ShorthandsOff}

\RequirePackage{geometry}
\RequirePackage{fancyhdr}
\RequirePackage{lastpage}
\RequirePackage{multicol}
\RequirePackage{array}
\RequirePackage{enumitem}
\RequirePackage{xcolor}
\RequirePackage{needspace}

\RequirePackage{amsmath,amssymb}
\RequirePackage{mathpazo}
\DeclareMathSizes{10.95}{10.95}{8}{5.5}
\begingroup\catcode`\_=\active \catcode`\^=\active
\gdef_#1{\sb{\scriptscriptstyle #1}}
\gdef^#1{\sp{\scriptscriptstyle\lycee@moinsepais #1}}
\endgroup
\newcommand\lycee@moins{\mathbin{%
  \setbox\z@\hbox{$\m@th\scriptscriptstyle\lycee@moinsorig$}%
  \dimen@=.55\wd\z@ \advance\dimen@-.5pt
  \hbox to.75\wd\z@{\hss$\m@th\scriptscriptstyle\vcenter{\hbox{%
    \tikz\draw[line width=.5pt,line cap=round](0,0)--(\the\dimen@,0);}}$\hss}}}
\begingroup\lccode`\~=`\- \lowercase{\endgroup\let~}\lycee@moins
\newcommand\lycee@moinsepais{\mathcode`\-="8000 }
\AtBeginDocument{\mathchardef\lycee@moinsorig=\mathcode`\-\relax}
\AtBeginDocument{\catcode`\_=12 \mathcode`\_="8000
  \catcode`\^=12 \mathcode`\^="8000 }
\newcommand\lycee@exposants{%
  \fontdimen13\textfont2=.48\fontdimen6\textfont2
  \fontdimen14\textfont2=.48\fontdimen6\textfont2
  \fontdimen15\textfont2=.48\fontdimen6\textfont2
  \fontdimen13\scriptfont2=.48\fontdimen6\scriptfont2
  \fontdimen14\scriptfont2=.48\fontdimen6\scriptfont2
  \fontdimen15\scriptfont2=.48\fontdimen6\scriptfont2}
\every@math@size\expandafter{\the\every@math@size\lycee@exposants}
\newlength{\retraitcalculs}
\setlength{\retraitcalculs}{2em}

\RequirePackage{tikz}
\usetikzlibrary{arrows.meta,calc,babel,shapes.misc}
\RequirePackage[most]{tcolorbox}
\tcbuselibrary{skins,breakable}

\definecolor{coulDef}{HTML}{1F4E79}
\definecolor{coulProp}{HTML}{7030A0}
\definecolor{coulRem}{HTML}{C55A11}
\definecolor{coulEx}{HTML}{2E7D32}
\definecolor{coulExo}{HTML}{B03A2E}
\definecolor{coulPart}{HTML}{1F4E79}
\definecolor{coulMeth}{HTML}{1B6E4F}
\definecolor{coulCourbe}{HTML}{0B5ED7}
\definecolor{coulCourbeB}{HTML}{C00000}
\definecolor{coulCourbeC}{HTML}{2E7D32}
\definecolor{coulGrille}{HTML}{8C8C8C}
\definecolor{coulRep}{HTML}{1B6E4F}
\definecolor{coulBar}{HTML}{2E7D32}

\newcommand{\R}{\mathbb{R}}
\newcommand{\Rs}{\mathbb{R}^{*}}
\renewcommand{\le}{\leqslant}
\renewcommand{\ge}{\geqslant}
\newcommand{\vect}[1]{\overrightarrow{#1}}
\newcommand{\norme}[1]{\left\lVert #1 \right\rVert}
\newcommand{\intervalle}[2]{\left[#1\,;#2\right]}
\RequirePackage{cancel}
\renewcommand{\CancelColor}{\color{coulExo}}
\newcommand{\fois}[1]{\times{\color{coulExo}#1}}
\newcommand{\barre}[1]{\cancel{#1}}

\tikzset{grille/.style={coulGrille,line width=0.4pt}}
\newcommand{\quadrillage}[4]{\draw[grille,step=1] (#1,#2) grid (#3,#4);}
\tikzset{
  croix/.style={cross out,draw,line width=0.6pt,inner sep=0pt,minimum size=4pt},
  croixdroite/.style={croix,rotate=45,line width=0.8pt},
}
\newcommand{\pt}[3]{\path (#1) node[croix]{} node[#2,font=\small,inner sep=2.5pt]{$#3$};}

\newcommand{\lycee@repere}[4]{%
  \draw[grille,step=1] (#1,#2) grid (#3,#4);
  \draw[-{Stealth[length=2mm]},thick] (#1,0)--({#3+0.4},0) node[below left=-1pt]{$x$};
  \draw[-{Stealth[length=2mm]},thick] (0,#2)--(0,{#4+0.4}) node[below left=-1pt]{$y$};
  \node[below left,font=\scriptsize,inner sep=1.5pt] at (0,0) {$0$};}
\newcommand{\lycee@gradx}[1]{\foreach \k/\l in {#1}{%
  \draw (\k,0.07)--(\k,-0.07) node[below,font=\scriptsize,inner sep=1.5pt]{$\l$};}}
\newcommand{\lycee@grady}[1]{\foreach \k/\l in {#1}{%
  \draw (0.07,\k)--(-0.07,\k) node[left,font=\scriptsize,inner sep=1.5pt]{$\l$};}}
\newcommand{\lycee@intff}[2]{\left[#1\,;#2\right]}
\newcommand{\lycee@intfo}[2]{\left[#1\,;#2\right[}
\newcommand{\lycee@intof}[2]{\left]#1\,;#2\right]}
\newcommand{\lycee@intoo}[2]{\left]#1\,;#2\right[}
\newcommand{\lycee@ens}[1]{\left\{#1\right\}}
\AtBeginDocument{%
  \providecommand{\repere}{\lycee@repere}%
  \providecommand{\gradx}{\lycee@gradx}%
  \providecommand{\grady}{\lycee@grady}%
  \providecommand{\Cf}{\mathcal{C}\sb{\scriptscriptstyle f}}%
  \providecommand{\Cg}{\mathcal{C}\sb{\scriptscriptstyle g}}%
  \providecommand{\intff}{\lycee@intff}%
  \providecommand{\intfo}{\lycee@intfo}%
  \providecommand{\intof}{\lycee@intof}%
  \providecommand{\intoo}{\lycee@intoo}%
  \providecommand{\ens}{\lycee@ens}}

\newlist{questions}{enumerate}{3}
\setlist[questions,1]{label=\arabic*\textdegree),leftmargin=*,itemsep=2pt,topsep=3pt}
\setlist[questions,2]{label=\alph*),leftmargin=*,itemsep=1pt,topsep=2pt}
\setlist[questions,3]{label=\roman*),leftmargin=*,itemsep=1pt,topsep=1pt}
\setlist[itemize]{label=\textbullet,leftmargin=*,itemsep=2pt,topsep=3pt}

\newcommand{\besoinplace}[1]{\par\penalty\@M\begingroup
  \dimen@=#1\relax\dimen@ii=\pagegoal\advance\dimen@ii-\pagetotal
  \ifdim\dimen@>\dimen@ii\ifdim\dimen@ii>\z@\vfil\fi\break\fi\endgroup}

\newcommand{\lycee@pied}{\thepage/\pageref{LastPage}}
\newcommand{\entete}[2]{%
  \pagestyle{fancy}\fancyhf{}%
  \fancyhead[L]{\footnotesize\color{coulPart}\textbf{#1}}%
  \fancyhead[R]{\footnotesize\color{coulPart}\textbf{#2}}%
  \fancyfoot[C]{\footnotesize\lycee@pied}%
  \renewcommand{\headrulewidth}{0.4pt}%
  \renewcommand{\headrule}{\hbox to\headwidth{%
    \color{coulPart!40}\leaders\hrule height \headrulewidth\hfill}}}

\RequirePackage{ccicons}
\newcommand{\lycee@licence}{{\scriptsize\color{gray}%
  \copyright\ Luc Destombe --- \ccbyncsa\ CC BY-NC-SA 4.0}}
\newif\iflycee@licence \lycee@licencetrue
\newcommand{\sanslicence}{\lycee@licencefalse}
\AtBeginDocument{\iflycee@licence\AddToHookNext{shipout/foreground}{%
  \put(0,0){\rlap{\kern\dimexpr1in+\hoffset+\oddsidemargin+\textwidth\relax
    \llap{\raisebox{-\dimexpr1in+\voffset+\topmargin+\headheight+\headsep
      +\textheight+\footskip\relax}{\lycee@licence}}}}}\fi}

\newcommand{\formulesserrees}{%
  \setlength{\abovedisplayskip}{3pt}\setlength{\belowdisplayskip}{3pt}%
  \setlength{\abovedisplayshortskip}{2pt}\setlength{\belowdisplayshortskip}{3pt}}

\setlength{\parindent}{0pt}

\RequirePackage[implicit=false,hidelinks,pdfencoding=auto,psdextra,bookmarksopen,
  bookmarksopenlevel=1,pdfcreator={},pdfproducer={}]{hyperref}
\RequirePackage{bookmark}
\hypersetup{pdfauthor={Luc Destombe},pdfkeywords={CC BY-NC-SA 4.0}}
\newcounter{lycee@signet}
\newcommand{\signet}[3][]{\stepcounter{lycee@signet}%
  \edef\lycee@dest{\if\relax\detokenize{#1}\relax
    signet.\arabic{lycee@signet}\else#1\fi}%
  \hypertarget{\lycee@dest}{}\bookmark[level=#2,dest=\lycee@dest]{#3}}
\pdfstringdefDisableCommands{%
  \def\R{R}\def\vect#1{#1}\def\Cf{Cf}\def\Cg{Cg}\def\fois#1{×#1}%
  \def\barre#1{#1}\def\intervalle#1#2{[#1;#2]}\def\intff#1#2{[#1;#2]}%
  \def\textsuperscript#1{#1}\def\dfrac#1#2{#1/#2}\def\frac#1#2{#1/#2}%
  \def\vec#1{vecteur #1}}
\geometry{top=2cm,bottom=1cm,left=1cm,right=1cm,headheight=14pt,footskip=0.6cm}

\RequirePackage{pgfplots}
\pgfplotsset{compat=1.18}
\RequirePackage{tkz-tab}

\newcounter{partiecours}
\renewcommand{\thepartiecours}{\Roman{partiecours}}
\newcommand{\partiecours}[1]{%
  \refstepcounter{partiecours}%
  \Needspace*{8\baselineskip}%
  \bigskip
  \noindent\begin{tcolorbox}[enhanced,colback=coulPart,colframe=coulPart,
      boxrule=0pt,arc=2pt,left=8pt,right=8pt,top=4pt,bottom=4pt,
      before skip=6pt,after skip=10pt]
    \color{white}\bfseries\large \thepartiecours{} --- #1\signet{1}{\thepartiecours{} --- #1}
  \end{tcolorbox}%
  \addcontentsline{toc}{section}{\thepartiecours{} --- #1}%
}

\newtcolorbox{boitecours}[3][]{%
  enhanced,breakable,
  colback=white,colframe=#2,coltitle=white,
  fonttitle=\bfseries,
  attach boxed title to top left={xshift=8pt,yshift=-\tcboxedtitleheight/2},
  boxed title style={colback=#2,arc=2pt,boxrule=0pt},
  boxrule=0.9pt,arc=2pt,
  left=8pt,right=8pt,top=8pt,bottom=6pt,
  title={#3},#1}

\newcounter{methode}
\newenvironment{definitionbox}[1][Définition]%
  {\begin{boitecours}[phantom={\signet{2}{#1}}]{coulDef}{#1}}{\end{boitecours}}
\newenvironment{proprietebox}[1][Propriété]%
  {\begin{boitecours}[phantom={\signet{2}{#1}}]{coulProp}{#1}}{\end{boitecours}}
\newenvironment{methodebox}[1][Méthode]{\stepcounter{methode}%
  \begin{boitecours}[phantom={\signet[methode-\arabic{methode}]{2}{#1}}]{coulMeth}{#1}}%
  {\end{boitecours}}

\newcommand{\etape}[1]{\par\smallskip\textbf{\color{coulMeth}#1}\par\nobreak\smallskip}
\newcommand{\enonce}[1]{\emph{#1}\par\smallskip}
\newcommand{\reponse}[1]{\par\smallskip
  {\footnotesize\itshape\color{black!55}Réponses : #1}}
\newcommand{\demo}[1]{\textbf{\color{coulMeth}Démonstration.} #1}
\newcommand{\afaire}[1]{%
  \par\smallskip
  \noindent{\small\color{coulExo}$\blacktriangleright$\ \textbf{À faire :}
    fiche d'exercices du chapitre, #1.}\par\smallskip}

\newenvironment{calculs}{\setlength{\jot}{5pt}\@fleqntrue\@mathmargin\retraitcalculs
  \setlength{\abovedisplayskip}{4pt}\setlength{\belowdisplayskip}{4pt}%
  \setlength{\abovedisplayshortskip}{4pt}\setlength{\belowdisplayshortskip}{4pt}%
  \csname alignat*\endcsname{2}}%
  {\csname endalignat*\endcsname}
\newcommand{\rem}[1]{\text{\small\itshape\color{black!60}#1}}
\newcommand{\explique}[2]{\par\smallskip\noindent
  \begin{minipage}[t]{0.57\linewidth}\raggedright #1\end{minipage}\hfill
  \begin{minipage}[t]{0.40\linewidth}\raggedright\leavevmode
    {\small\itshape\color{black!60}#2}\end{minipage}\par}
\newcommand{\cas}[1]{\par\smallskip\noindent\textbf{\color{coulExo}#1}\nobreak}
\newcommand{\calc}[1]{\par\smallskip\textbf{\color{coulExo}#1}\par\nobreak}

\newtcolorbox{remarquebox}[1][Remarque]{%
  enhanced,breakable,colback=white,colframe=coulRem,
  boxrule=0pt,leftrule=2.5pt,arc=0pt,outer arc=0pt,
  left=8pt,right=6pt,top=5pt,bottom=5pt,
  before upper={\textbf{\color{coulRem}#1}\par\smallskip}}

\newtcolorbox{exemplebox}[1][Exemple]{%
  enhanced,breakable,colback=white,colframe=coulEx,
  boxrule=0pt,leftrule=2.5pt,arc=0pt,outer arc=0pt,
  left=8pt,right=6pt,top=5pt,bottom=5pt,
  before upper={\textbf{\color{coulEx}#1}\par\smallskip}}

\pgfplotsset{
  repere/.style={
    axis lines=middle,
    axis line style={-{Stealth[length=2.4mm]},thick},
    every axis x label/.style={at={(ticklabel* cs:1.0)},anchor=west},
    every axis y label/.style={at={(ticklabel* cs:1.0)},anchor=south},
    label style={font=\footnotesize},
    tick label style={font=\scriptsize},
    grid=both,
    major grid style={grille},
    minor tick num=0,
    tick align=inside,
    clip mode=individual,
  },
}
\tikzset{
  courbe/.style={coulCourbe,line width=1.1pt,smooth},
  courbeB/.style={coulCourbeB,line width=1.1pt,smooth},
  courbeC/.style={coulCourbeC,line width=1.1pt,smooth},
}

\newlist{etapes}{enumerate}{1}
\setlist[etapes]{label=\textbf{\arabic*.},leftmargin=*,itemsep=1pt,topsep=2pt}

\newlist{expressions}{enumerate}{1}
\setlist[expressions]{label={\Alph*\ $=$},leftmargin=*,itemsep=3pt,topsep=3pt}

\newcounter{exo}
\renewcommand{\theexo}{\ifnum\value{exo}<10 0\fi\arabic{exo}}
\newtcolorbox{exercice}[1][]{%
  enhanced,breakable,
  colback=white,colframe=coulExo,
  boxrule=0.9pt,arc=2pt,
  left=8pt,right=8pt,top=8pt,bottom=6pt,
  coltitle=white,fonttitle=\bfseries,
  attach boxed title to top left={xshift=8pt,yshift=-\tcboxedtitleheight/2},
  boxed title style={colback=coulExo,arc=2pt,boxrule=0pt},
  phantom={\signet{2}{Exercice \arabic{exo}}},
  title={Exercice \theexo\ifx\relax#1\relax\else{} --- #1\fi}}
\newenvironment{exo}[1][\relax]{\refstepcounter{exo}\begin{exercice}[#1]}{\end{exercice}}

  \RequirePackage{comment}
  \excludecomment{correction}
  \renewcommand{\reponse}[1]{}
  \newcommand{\autableau}{\hfill{\footnotesize\itshape\color{black!45}(au tableau)}}
  \newcommand{\etapeexemples}{%
    \par\smallskip\textbf{\color{coulMeth}Exemples}\autableau\par\nobreak\smallskip}
  \newcommand{\etapetableau}[1]{%
    \par\smallskip\textbf{\color{coulMeth}#1}\autableau\par\nobreak\smallskip}
  \newtcolorbox{exempletableau}[1][Exemple]{%
    enhanced,breakable,colback=white,colframe=coulEx,
    boxrule=0pt,leftrule=2.5pt,arc=0pt,outer arc=0pt,
    left=8pt,right=6pt,top=4pt,bottom=4pt,
    before skip=5pt,after skip=5pt,
    before upper={\textbf{\color{coulEx}#1}\autableau\par\smallskip}}

\newcommand{\coursEntete}{}
\newcommand{\coursNiveau}{}
\newcommand{\coursTitre}{}
\newcommand{\coursSurTitre}{}
\newcommand{\coursSousTitre}{}

\newcommand{\enteteCours}[1]{\renewcommand{\coursEntete}{#1}}
\newcommand{\chapitrecours}[2]{\enteteCours{Chapitre #1 --- #2}}
\newcommand{\niveaucours}[1]{\renewcommand{\coursNiveau}{#1}}
\newcommand{\titrecours}[3][]{%
  \renewcommand{\coursTitre}{#2}%
  \renewcommand{\coursSurTitre}{#1}%
  \renewcommand{\coursSousTitre}{#3}}

\pagestyle{fancy}
\fancyhf{}
\fancyhead[L]{\footnotesize\color{coulPart}\textbf{\coursEntete}}
\fancyhead[R]{\footnotesize\color{coulPart}\textbf{\coursNiveau}}
\fancyfoot[C]{\footnotesize\thepage}
\renewcommand{\headrulewidth}{0.4pt}
\renewcommand{\headrule}{\hbox to\headwidth{\color{coulPart!40}\leaders\hrule height \headrulewidth\hfill}}

\setlength{\parskip}{2pt}

\AtBeginDocument{%
  \begin{center}
    \begin{tcolorbox}[enhanced,width=0.72\linewidth,colback=white,colframe=coulPart,
        boxrule=1.6pt,arc=3pt,halign=center,top=10pt,bottom=10pt]
      \ifx\coursSurTitre\empty
        {\Huge\bfseries\color{coulPart} \coursTitre}\\[4pt]
      \else
        {\Huge\bfseries\color{coulPart} \coursTitre}\\[3pt]
        {\Large\bfseries\color{coulPart} \coursSurTitre}\\[5pt]
      \fi
      {\normalsize \coursSousTitre}%
      \\[2pt]{\small\itshape Version élève}
    \end{tcolorbox}
  \end{center}%
}
\makeatother

\chapitrecours{2}{Fonctions affines et pourcentages}
\niveaucours{1\textsuperscript{re} mathématiques spécifiques}
\titrecours{FONCTIONS AFFINES ET POURCENTAGES}{Cours --- Première, mathématiques spécifiques}

\begin{document}

\begin{remarquebox}[Comment utiliser ce cours]
Chaque partie commence par ce qu'il faut \textbf{savoir} (définitions, propriétés), puis
vient ce qu'il faut \textbf{savoir faire} : une méthode, avec des
\textbf{exemples} faits en classe, puis des exercices
\og\textbf{À vous de jouer}\fg{} à chercher seul. Tout le chapitre se fait
\textbf{sans calculatrice}.

\smallskip
Sur cette feuille, seuls les \textbf{énoncés} des exemples figurent : la résolution,
signalée par la mention \emph{(au tableau)}, est à rédiger dans le cahier.

\end{remarquebox}

\partiecours{Appliquer un pourcentage}

\begin{definitionbox}[Définition 1 --- Pourcentage]
$t\,\%$ est une autre façon d'écrire le nombre $\dfrac{t}{100}$. Par exemple :
\[
  35\,\%=\dfrac{35}{100}=0{,}35 \qquad\quad 6\,\%=\dfrac{6}{100}=0{,}06 \qquad\quad
  120\,\%=\dfrac{120}{100}=1{,}2
\]
Prendre $t\,\%$ d'une quantité, c'est la \textbf{multiplier} par $\dfrac{t}{100}$.
\end{definitionbox}

\begin{remarquebox}[Pour calculer de tête]
\begin{center}
\renewcommand{\arraystretch}{1.4}
\begin{tabular}{|c|c|c|c|c|c|}
\hline
$1\,\%$ & $10\,\%$ & $5\,\%$ & $25\,\%$ & $50\,\%$ & $75\,\%$\\
\hline
diviser par $100$ & diviser par $10$ & moitié de $10\,\%$ & le quart &
la moitié & trois quarts\\
\hline
\end{tabular}
\end{center}
\end{remarquebox}

\begin{methodebox}[Méthode 1 --- Calculer $t\,\%$ d'une quantité]
\etapeexemples
\begin{questions}[label=\alph*)]
  \item Calculer $40\,\%$ de $120$.
  \item Une salle de concert a $1\,200$ places ; $15\,\%$ sont réservées aux abonnés.
        Combien de places sont réservées ?
  \item Dans une ville, $30$ vélos en libre-service représentent $5\,\%$ du parc.
        Combien le parc compte-t-il de vélos ?
\end{questions}

\etape{À vous de jouer}
\begin{questions}
  \item Calculer $25\,\%$ de $360$ ; $10\,\%$ de $48$ ; $75\,\%$ de $60$.
  \item Dans une classe de $30$ élèves, $40\,\%$ font du sport en club. Combien
        d'élèves cela fait-il ?
  \item $18$ personnes représentent $25\,\%$ d'un groupe. Combien le groupe
        compte-t-il de personnes ?
\end{questions}

\end{methodebox}

\begin{methodebox}[Méthode 2 --- Augmenter ou diminuer une quantité de $t\,\%$]
\etapeexemples
\begin{questions}[label=\alph*)]
  \item Une console à $300$~€ augmente de $10\,\%$. Calculer son nouveau prix.
  \item Un manteau à $120$~€ est soldé à $-25\,\%$. Calculer son prix soldé.
\end{questions}

\etape{À vous de jouer}
\begin{questions}
  \item Un abonnement de $400$~€ augmente de $3\,\%$. Calculer le nouveau prix.
  \item Une commune de $8\,000$ habitants perd $5\,\%$ de sa population. Combien
        d'habitants compte-t-elle ensuite ?
  \item Un sac à $60$~€ est soldé à $-30\,\%$. Calculer son prix soldé.
\end{questions}

\end{methodebox}

\afaire{exercices 1 à 5}

\partiecours{Calculer un pourcentage}

\begin{definitionbox}[Définition 2 --- Proportion]
Dans un ensemble de $N$ individus (le « total »), une partie en compte $n$. La
\textbf{proportion} de cette partie est
\[
  \boxed{p=\dfrac{n}{N}}
\]
C'est un nombre entre $0$ et $1$, qu'on écrit souvent en pourcentage.
\end{definitionbox}

\begin{methodebox}[Méthode 3 --- Calculer une proportion en pourcentage]
\etapeexemples
\begin{questions}[label=\alph*)]
  \item Sur $20$ élèves, $9$ viennent au lycée à vélo. Quel pourcentage des élèves
        cela représente-t-il ?
  \item Un spectacle réunit $300$ spectateurs, dont $54$ enfants. Quelle est la
        proportion d'enfants, en pourcentage ?
\end{questions}

\etape{À vous de jouer}
Écrire chaque proportion en pourcentage :
\[
  \dfrac{7}{20} \qquad\qquad \dfrac{1}{4} \qquad\qquad \dfrac{36}{200}
  \qquad\qquad \text{$18$ reçus sur $60$ candidats}
\]

\end{methodebox}

\begin{definitionbox}[Définition 3 --- Taux d'évolution]
Une grandeur passe d'une valeur de départ $V_{0}$ à une valeur d'arrivée $V_{1}$. Son
\textbf{taux d'évolution} est
\[
  \boxed{t=\dfrac{V_{1}-V_{0}}{V_{0}}}
\]
\begin{itemize}
  \item $t>0$ : c'est une \textbf{hausse} ; $t<0$ : c'est une \textbf{baisse} ;
  \item on divise toujours par la valeur de \textbf{départ}.
\end{itemize}
\end{definitionbox}

\begin{methodebox}[Méthode 4 --- Calculer un taux d'évolution]
\etapeexemples
\begin{questions}[label=\alph*)]
  \item Un prix passe de $60$~€ à $69$~€. Calculer le taux d'évolution.
  \item Une récolte passe de $500$ à $470$ tonnes. Calculer le taux d'évolution.
  \item Une association passe de $600$ à $530$ membres. Justifier qu'elle a perdu plus
        de $10\,\%$ de ses membres.
\end{questions}

\etape{À vous de jouer}
Calculer le taux d'évolution de chaque grandeur :
\begin{questions}
  \item un salaire qui passe de $2\,000$~€ à $2\,080$~€ ;
  \item une facture qui passe de $40$~€ à $30$~€ ;
  \item une production qui passe de $120$ à $150$ tonnes ;
  \item une ville passe de $4\,000$ à $4\,300$ habitants : a-t-elle gagné plus de
        $5\,\%$ d'habitants ?
\end{questions}

\end{methodebox}

\afaire{exercices 6 à 10}

\partiecours{Coefficient multiplicateur}

\begin{exemplebox}[Pour commencer]
Augmenter $150$ de $20\,\%$, c'est ajouter $0{,}2\times 150$ :
\[
  150+0{,}2\times 150=150\times(1+0{,}2) \qquad\text{soit}\qquad 150\times 1{,}2=180
\]
Le résultat s'obtient en \textbf{une seule multiplication}.
\end{exemplebox}

\begin{proprietebox}[Propriété 1 --- Coefficient multiplicateur]
\begin{itemize}
  \item Augmenter de $t\,\%$, c'est multiplier par $\boxed{1+\dfrac{t}{100}}$.
  \item Diminuer de $t\,\%$, c'est multiplier par $\boxed{1-\dfrac{t}{100}}$.
\end{itemize}
Ce nombre est le \textbf{coefficient multiplicateur} $C_{M}$ : $V_{1}=V_{0}\times C_{M}$.
\begin{itemize}
  \item $C_{M}>1$ : c'est une hausse ; $C_{M}<1$ : c'est une baisse.
  \item On lit le taux d'évolution sur $C_{M}$ : $t=C_{M}-1$.
\end{itemize}
\end{proprietebox}

\begin{center}
\renewcommand{\arraystretch}{1.4}
\begin{tabular}{|l|c|c|c|c|c|c|}
\hline
Évolution & $+30\,\%$ & $+4\,\%$ & $+100\,\%$ & $-20\,\%$ & $-6\,\%$ & $-50\,\%$\\
\hline
Coefficient multiplicateur & $1{,}3$ & $1{,}04$ & $2$ & $0{,}8$ & $0{,}94$ & $0{,}5$\\
\hline
\end{tabular}
\end{center}

\begin{methodebox}[Méthode 5 --- Appliquer une évolution avec le coefficient multiplicateur]
\etapeexemples
\begin{questions}[label=\alph*)]
  \item Un téléphone à $400$~€ baisse de $15\,\%$. Quel calcul donne son nouveau prix ?
  \item Une colonie de $2\,000$ abeilles augmente de $40\,\%$ au printemps. Combien
        d'abeilles compte-t-elle ensuite ?
  \item Un prix est multiplié par $0{,}94$, un autre par $1{,}15$. Quelles évolutions
        ont-ils subies ?
  \item Un billet passe de $80$~€ à $100$~€. Retrouver le taux d'évolution avec le
        coefficient multiplicateur.
\end{questions}

\etape{À vous de jouer}
\begin{questions}
  \item Donner le coefficient multiplicateur de : $+8\,\%$ ; $-40\,\%$ ; $+200\,\%$.
  \item Donner l'évolution associée à : $\times 1{,}05$ ; $\times\dfrac{1}{2}$ ;
        $\times 2$.
  \item Un salaire de $1\,800$~€ augmente de $5\,\%$ ; une facture de $60$~€ baisse de
        $25\,\%$. Calculer les nouveaux montants.
  \item Un prix passe de $50$~€ à $40$~€. Calculer $C_{M}$, puis le taux d'évolution.
\end{questions}

\end{methodebox}

\afaire{exercices 11 à 15}

\partiecours{Évolutions successives, évolution réciproque}

\begin{exemplebox}[Pour commencer]
Un jean à $80$~€ baisse de $20\,\%$, puis augmente de $20\,\%$ :
\[
  80\times 0{,}8=64 \qquad\text{puis}\qquad 64\times 1{,}2=76{,}80
\]
Il ne revient pas à $80$~€ : les pourcentages \textbf{ne s'ajoutent pas}.
\end{exemplebox}

\begin{proprietebox}[Propriété 2 --- Évolutions successives]
Le coefficient multiplicateur global de plusieurs évolutions successives est le
\textbf{produit} de leurs coefficients multiplicateurs :
\[
  \boxed{C_{M\,\text{global}}=C_{M1}\times C_{M2}\times\cdots}
\]
Une même évolution répétée $n$ fois a pour coefficient global $C_{M}^{n}$ (suite
géométrique de raison $C_{M}$, chapitre~1).
\end{proprietebox}

\begin{methodebox}[Méthode 6 --- Calculer un taux d'évolution global]
\etapeexemples
Calculer le taux d'évolution global dans chaque cas.
\begin{questions}[label=\alph*)]
  \item Un prix augmente de $20\,\%$, puis baisse de $20\,\%$.
  \item Le chiffre d'affaires d'un magasin augmente de $10\,\%$ en 2025, puis de
        $30\,\%$ en 2026.
  \item Un loyer augmente de $4\,\%$ par an pendant $3$ ans. On donne
        $1{,}04^{3}\approx 1{,}12$.
\end{questions}

\etape{À vous de jouer}
Calculer le taux d'évolution global dans chaque cas :
\begin{questions}
  \item une hausse de $50\,\%$, puis une baisse de $40\,\%$ ;
  \item deux baisses successives de $10\,\%$ ;
  \item une hausse de $20\,\%$ par an pendant $2$ ans.
\end{questions}

\end{methodebox}

\begin{definitionbox}[Définition 4 --- Évolution réciproque]
L'\textbf{évolution réciproque} d'une évolution est celle qui ramène à la valeur de
départ. Deux évolutions successives se \textbf{compensent} quand le produit de leurs
coefficients multiplicateurs vaut $1$ :
\[
  \boxed{C_{M1}\times C_{M2}=1}
\]
\end{definitionbox}

\begin{remarquebox}[Cas à connaître]
\begin{center}
\renewcommand{\arraystretch}{1.4}
\begin{tabular}{|l|c|c|c|}
\hline
Évolution & $+25\,\%$ & $+100\,\%$ (doubler) & $+300\,\%$ (quadrupler)\\
\hline
Compensée par & $-20\,\%$ & $-50\,\%$ (moitié) & $-75\,\%$ (quart)\\
\hline
Vérification & $1{,}25\times 0{,}8=1$ & $2\times 0{,}5=1$ & $4\times 0{,}25=1$\\
\hline
\end{tabular}
\end{center}
Une hausse de $t\,\%$ n'est \textbf{jamais} compensée par une baisse de $t\,\%$.
\end{remarquebox}

\begin{methodebox}[Méthode 7 --- Vérifier qu'une évolution en compense une autre]
\etapeexemples
\begin{questions}[label=\alph*)]
  \item Un prix double. Une baisse de $50\,\%$ le ramène-t-elle à sa valeur de départ ?
  \item Une hausse de $20\,\%$ est-elle compensée par une baisse de $20\,\%$ ?
  \item Une action perd $20\,\%$ de sa valeur. Pour la retrouver, il faut une hausse
        de : $20\,\%$ ; $25\,\%$ ; $40\,\%$ ?
\end{questions}

\etape{À vous de jouer}
\begin{questions}
  \item Une baisse de $50\,\%$ est-elle compensée par une hausse de $100\,\%$ ?
  \item Une hausse de $10\,\%$ est-elle compensée par une baisse de $10\,\%$ ?
  \item Un prix est multiplié par $4$. Pour revenir au prix de départ, il faut une
        baisse de : $75\,\%$ ; $300\,\%$ ; $25\,\%$ ?
\end{questions}

\end{methodebox}

\afaire{exercices 16 à 23}

\partiecours{Fonctions affines : image, antécédent}

\begin{definitionbox}[Définition 5 --- Fonction affine]
Une fonction $f$ est \textbf{affine} quand il existe deux nombres $a$ et $b$ tels que,
pour tout réel $x$ :
\[
  \boxed{f(x)=ax+b}
\]
$a$ est le \textbf{coefficient directeur}, $b$ l'\textbf{ordonnée à l'origine}.
\begin{itemize}
  \item Si $b=0$, $f(x)=ax$ : $f$ est \textbf{linéaire} (proportionnalité).
  \item Si $a=0$, $f(x)=b$ : $f$ est \textbf{constante}.
\end{itemize}
\end{definitionbox}

\begin{exemplebox}[Exemples]
\begin{itemize}
  \item Une course en VTC coûte $5$~€ de prise en charge, puis $2$~€ par kilomètre.
        Pour $x$~km, on paie $f(x)=2x+5$ : $a=2$ et $b=5$.
  \item Baisser un prix $x$ de $10\,\%$ donne $h(x)=0{,}9x$ : $h$ est linéaire.
\end{itemize}
\end{exemplebox}

\begin{remarquebox}[Image et antécédent (rappel de seconde)]
Si $f(x)=y$, on dit que $y$ est l'\textbf{image} de $x$ et que $x$ est un
\textbf{antécédent} de $y$. Pour une fonction affine non constante, chaque nombre a
un seul antécédent.
\end{remarquebox}

\begin{methodebox}[Méthode 8 --- Calculer une image]
\etapeexemples
\begin{questions}[label=\alph*)]
  \item Soit $f(x)=-3x+8$. Calculer $f(2)$ et $f(-5)$.
  \item VTC : calculer le prix d'une course de $12$~km.
\end{questions}

\etape{À vous de jouer}
\begin{questions}
  \item Soit $f(x)=4x-7$. Calculer $f(0)$, $f(3)$ et $f(-2)$.
  \item Soit $h(x)=0{,}5x+3$. Calculer $h(10)$ et $h(-4)$.
  \item VTC : calculer le prix d'une course de $8$~km.
\end{questions}

\end{methodebox}

\begin{methodebox}[Méthode 9 --- Calculer un antécédent]
\etapeexemples
\begin{questions}[label=\alph*)]
  \item Soit $f(x)=4x-7$. Calculer l'antécédent de $9$.
  \item VTC : une course a coûté $35$~€. Quelle distance a-t-on parcourue ?
\end{questions}

\etape{À vous de jouer}
\begin{questions}
  \item Soit $f(x)=-3x+8$. Calculer l'antécédent de $14$, puis celui de $2$.
  \item VTC : une course a coûté $45$~€. Quelle distance a-t-on parcourue ?
  \item Soit $h(x)=5x+2$. Résoudre $h(x)=-8$.
\end{questions}

\end{methodebox}

\afaire{exercices 24 à 27}

\partiecours{Trouver une fonction affine}

\begin{proprietebox}[Propriété 3 --- Coefficient directeur à partir de deux points]
Si la droite qui représente $f(x)=ax+b$ passe par $A(x_{A}\,;y_{A})$ et
$B(x_{B}\,;y_{B})$, alors :
\[
  \boxed{a=\dfrac{y_{B}-y_{A}}{x_{B}-x_{A}}}
  \qquad\text{(différence des ordonnées sur différence des abscisses)}
\]
Quand $x$ augmente de $1$, $f(x)$ augmente de $a$ (diminue si $a<0$) ; et $b=f(0)$ est
la valeur de départ.
\end{proprietebox}

\begin{methodebox}[Méthode 10 --- Déterminer une fonction affine]
\etapeexemples
\begin{questions}[label=\alph*)]
  \item Un abonnement de vidéo coûte $9$~€ par mois, plus $3$~€ par film loué.
        Exprimer le prix mensuel $p(x)$ pour $x$ films loués.
  \item La droite qui représente $f$ passe par $A(1\,;4)$ et $B(4\,;13)$. Déterminer
        $f(x)$.
  \item La droite qui représente $g$ passe par $A(-2\,;5)$ et $B(-4\,;9)$. Déterminer
        $g(x)$.
  \item Une facture d'électricité vaut $70$~€ pour $200$~kWh et $100$~€ pour
        $300$~kWh. On la modélise par une fonction affine $F$ de la consommation $x$.
        Déterminer $F(x)$, puis interpréter $a$ et $b$.
\end{questions}

\etape{À vous de jouer}
\begin{questions}
  \item Un club de danse coûte $40$~€ d'inscription, puis $20$~€ par mois. Exprimer le
        coût $C(x)$ pour $x$ mois.
  \item La droite qui représente $g$ passe par $A(2\,;9)$ et $B(5\,;3)$. Déterminer
        $g(x)$.
  \item La droite qui représente $h$ passe par $A(-1\,;-4)$ et $B(3\,;8)$. Déterminer
        $h(x)$.
\end{questions}

\end{methodebox}

\afaire{exercices 28 à 31}

\partiecours{Tracer une fonction affine}

\begin{proprietebox}[Propriété 4 --- Représentation graphique et variations]
La représentation graphique de $f(x)=ax+b$ est une \textbf{droite}, d'\textbf{équation
réduite} $\boxed{y=ax+b}$ : le point $M(x\,;y)$ est sur la droite quand $y=ax+b$.
\begin{itemize}
  \item elle coupe l'axe des ordonnées au point $(0\,;b)$ ;
  \item pour un déplacement horizontal de $1$, le déplacement vertical vaut $a$ (on
        monte si $a>0$, on descend si $a<0$).
\end{itemize}
Sens de variation : $a>0$, $f$ est \textbf{croissante} ; $a<0$, $f$ est
\textbf{décroissante} ; $a=0$, $f$ est \textbf{constante}.
\end{proprietebox}

\begin{remarquebox}[Lien avec les suites arithmétiques]
Les points d'une suite arithmétique $u_{n}=u_{0}+rn$ sont sur la droite qui représente
$f(x)=rx+u_{0}$ : dans les deux cas, on parle de \textbf{croissance linéaire}.
\end{remarquebox}

\begin{methodebox}[Méthode 11 --- Tracer la droite d'une fonction affine]
\etapeexemples
Tracer les droites qui représentent $f(x)=3x-2$ et $g(x)=-0{,}5x+3$.

\etape{À vous de jouer}
Tracer dans un repère les droites qui représentent $h(x)=x-2$ ; $k(x)=-2x+3$ et
$m(x)=\dfrac{1}{2}x+1$.

\end{methodebox}

\begin{proprietebox}[Propriété 5 --- Fonctions linéaires et proportionnalité]
Une fonction \textbf{linéaire} $f(x)=ax$ (cas $b=0$) traduit une situation de
\textbf{proportionnalité} : $a$ est le coefficient de proportionnalité, la quantité
pour une unité de $x$.
\begin{itemize}
  \item Sa droite passe par l'\textbf{origine} du repère.
  \item Appliquer une évolution de $t\,\%$ est une fonction linéaire :
        $x\mapsto C_{M}\times x$.
\end{itemize}
\end{proprietebox}

\begin{methodebox}[Méthode 12 --- Utiliser une fonction linéaire (proportionnalité)]
\etapeexemples
\begin{questions}[label=\alph*)]
  \item Des céréales contiennent $8$~g de fibres pour $100$~g. Quelle masse de fibres
        contient un paquet de $350$~g ?
  \item Un mile vaut environ $1{,}6$~km. Convertir $50$ miles en kilomètres, puis
        $24$~km en miles.
  \item Parmi $f(x)=-3x$ ; $g(x)=3x-1$ et $h(x)=0{,}4x$, lesquelles sont linéaires ?
\end{questions}

\etape{À vous de jouer}
\begin{questions}
  \item Une imprimante imprime $20$ pages par minute. Exprimer le nombre $P(t)$ de pages
        imprimées en $t$ minutes ; calculer $P(4)$, puis le temps pour imprimer
        $300$ pages.
  \item En soldes, un prix $x$ baisse de $15\,\%$. Exprimer le prix soldé $R(x)$ ;
        calculer $R(60)$, puis le prix de départ d'un article soldé à $34$~€.
  \item Tracer la droite qui représente $k(x)=3x$.
\end{questions}

\end{methodebox}

\afaire{exercices 32 à 35}

\partiecours{Lire l'équation d'une droite}

\begin{proprietebox}[Propriété 6 --- Lire $a$ et $b$ sur la droite]
Sur la droite d'équation réduite $y=ax+b$ :
\begin{itemize}
  \item $b$ est l'ordonnée du point où la droite coupe l'axe des ordonnées ;
  \item entre deux points de la droite,
        $\boxed{a=\dfrac{\text{déplacement vertical}}{\text{déplacement horizontal}}}$,
        le déplacement vertical étant négatif quand on descend (c'est
        $\dfrac{y_{B}-y_{A}}{x_{B}-x_{A}}$ lu sur le quadrillage).
\end{itemize}
\end{proprietebox}

\begin{methodebox}[Méthode 13 --- Lire l'équation réduite d'une droite]
\etapeexemples
\noindent
\begin{minipage}[c]{0.58\linewidth}\raggedright
Les droites $(d_{1})$ et $(d_{2})$ représentent deux fonctions affines $u$ et $v$.
Donner l'équation réduite de chaque droite, puis le sens de variation de $u$ et de $v$.
\end{minipage}\hfill
\begin{minipage}[c]{0.38\linewidth}\centering
\begin{tikzpicture}
\begin{axis}[repere,
    x=0.42cm,y=0.42cm,
    xmin=-1.5,xmax=6.5,ymin=-3.5,ymax=4.5,
    xtick={-1,...,6},ytick={-3,...,4},
    xlabel={$x$},ylabel={$y$}]
  \addplot[coulCourbe,line width=1.1pt,domain=-0.6:3.2]{-2*x+3};
  \addplot[coulCourbeB,line width=1.1pt,domain=-1.5:6.5]{x/3-1};
  \node[coulCourbe,right] at (axis cs:-0.5,4) {$(d_{1})$};
  \node[coulCourbeB,above] at (axis cs:5.4,1) {$(d_{2})$};
\end{axis}
\end{tikzpicture}
\end{minipage}

\etape{À vous de jouer}
Une droite passe par $(0\,;-3)$ et $(1\,;1)$ ; une autre par $(0\,;2)$ et $(2\,;1)$.
Donner l'équation réduite de chaque droite et le sens de variation de la fonction
qu'elle représente.

\end{methodebox}

\afaire{exercices 36 à 38}

\partiecours{Signe d'une fonction affine}

\begin{proprietebox}[Propriété 7 --- Signe de $ax+b$ (avec $a\ne 0$)]
$ax+b$ s'annule pour une seule valeur, $x=-\dfrac{b}{a}$ (solution de $ax+b=0$).
\begin{itemize}
  \item Si $a>0$ (la droite monte), $ax+b$ est \textbf{négatif avant} cette valeur et
        \textbf{positif après}.
  \item Si $a<0$ (la droite descend), $ax+b$ est \textbf{positif avant} et
        \textbf{négatif après}.
\end{itemize}
À retenir : dans le tableau de signes, à droite du zéro, on écrit le \textbf{signe de
$a$}.
\end{proprietebox}

\begin{methodebox}[Méthode 14 --- Dresser le tableau de signes d'une fonction affine]
\etapeexemples
\begin{questions}[label=\alph*)]
  \item Dresser le tableau de signes de $f(x)=3x-12$.
  \item Dresser le tableau de signes de $g(x)=-5x+2$.
  \item Une entreprise vend $x$ objets ; son bénéfice, en euros, est $B(x)=4x-60$. À
        partir de combien d'objets vendus le bénéfice est-il positif ?
\end{questions}

\etape{À vous de jouer}
Dresser le tableau de signes de chaque fonction :
\[
  f(x)=2x+8 \qquad\qquad g(x)=-x+3 \qquad\qquad h(x)=0{,}5x-3
\]

\end{methodebox}

\afaire{exercices 39 à 41}

\partiecours{Comparer deux fonctions affines}

\begin{remarquebox}[Rappel --- Résoudre une équation $ax+b=cx+d$]
On regroupe les termes en $x$ d'un côté et les nombres de l'autre, puis on divise par le
coefficient de $x$ :
\begin{calculs}
7x+2 &= 3x+14 &&\\
4x &= 12 &\qquad& \rem{on retire $3x$ et $2$}\\
x &= 3 &\qquad& \rem{on divise par $4$}
\end{calculs}
\end{remarquebox}

\begin{methodebox}[Méthode 15 --- Comparer deux offres]
\etapeexemples
\noindent
\begin{minipage}[c]{0.58\linewidth}\raggedright
Une agence loue une camionnette à la journée. Tarif A : $20$~€, plus $0{,}50$~€ par
km ; tarif B : $80$~€, plus $0{,}20$~€ par km. On note $f(x)$ et $g(x)$ les prix pour
$x$~km. Quel tarif choisir selon la distance ? Le graphique ci-contre représente les
deux tarifs.
\end{minipage}\hfill
\begin{minipage}[c]{0.38\linewidth}\centering
\begin{tikzpicture}
\begin{axis}[repere,
    x=0.016cm,y=0.016cm,
    xmin=0,xmax=330,ymin=0,ymax=200,
    xtick={50,100,...,300},ytick={20,40,...,180},
    xticklabels={,100,,200,,300},yticklabels={,40,,80,,120,,160,},
    xlabel={km},ylabel={€}]
  \addplot[coulCourbe,line width=1.1pt,domain=0:320]{0.5*x+20};
  \addplot[coulCourbeB,line width=1.1pt,domain=0:320]{0.2*x+80};
  \addplot[only marks,mark=x,mark options={line width=0.6pt},mark size=2pt]
    coordinates {(200,120)};
  \node[coulCourbe,left] at (axis cs:280,170) {A};
  \node[coulCourbeB,below] at (axis cs:300,140) {B};
\end{axis}
\end{tikzpicture}
\end{minipage}

\etape{À vous de jouer}
\begin{questions}
  \item Résoudre $4x+3=x+12$ ; puis $6x-5=2x+11$ ; puis $2x+5=-4x+8$.
  \item Salle A : $20$~€ d'inscription et $30$~€ par mois ; salle B : $100$~€
        d'inscription et $20$~€ par mois. À partir de combien de mois la salle B
        est-elle la moins chère ?
\end{questions}

\end{methodebox}

\afaire{exercices 42 à 44, puis les exercices type épreuve 45 à 47}

\end{document}
